\chapter{Assemblers}

\section{The job of an assembler}

An \textbf{assembler} translates assembly language to machine code. Since assembly is (almost)
one-to-one with machine code, the interesting parts are:
\begin{enumerate}
  \item \textbf{Analysis}: break the text into tokens, check each line is well formed, and
        collect label definitions.
  \item \textbf{Synthesis}: produce the 32-bit word for each instruction, resolving label
        references.
\end{enumerate}
Output is binary written to standard output (four bytes per word, big-endian); errors are
reported on standard error, and on any error no (partial) output should be considered valid.

\section{Scanning the input}

The course supplied a scanner (in C++ and Racket) that splits each line into tokens with
kinds such as \code{ID}, \code{LABEL} (\code{foo:}), \code{REG} (\code{\$12}),
\code{INT} (\code{-17}), \code{HEXINT} (\code{0x1f}), \code{COMMA}, \code{LPAREN},
\code{RPAREN}, \code{DOTID} (\code{.word}) and \code{NEWLINE}. How scanners work in general is
the subject of Chapter~7; here we simply consume the token stream.

\section{The forward-reference problem}

A label may be used before it is defined:
\begin{lstlisting}[style=mips]
    beq $1, $0, done     ; 'done' not seen yet!
    add $3, $3, $1
done:
    jr $31
\end{lstlisting}
When we reach the \code{beq} we do not know \code{done}'s address. Two solutions:
\begin{itemize}
  \item \textbf{Two passes}: first pass computes all label addresses; second pass encodes.
  \item \textbf{One pass with backpatching}: emit a placeholder and record the location; patch
        it when the label is defined (requires buffering output).
\end{itemize}
CS 241 uses the two-pass approach.

\section{Two-pass assembly}

\subsection{Pass 1: analysis and the symbol table}
Maintain a \emph{location counter} \code{PC}, starting at 0.
\begin{lstlisting}[style=cpp]
std::map<std::string, uint32_t> symtab;
uint32_t pc = 0;
for (auto &line : lines) {
    size_t k = 0;
    while (k < line.size() && line[k].kind == LABEL) {
        std::string name = labelName(line[k]);      // strip ':'
        if (symtab.count(name)) error("duplicate label " + name);
        symtab[name] = pc;
        ++k;
    }
    if (k < line.size()) {               // an instruction or .word
        checkSyntax(line, k);            // operand kinds, counts, commas
        pc += 4;
    }
}
\end{lstlisting}
Key points: a label alone on a line (or followed only by a comment) does not advance
\code{PC}; labels mark the \emph{next} instruction even if that is several lines later.

\subsection{Pass 2: synthesis}
Walk the lines again with a fresh \code{PC}. For each instruction, look up the opcode/funct,
check operand ranges, substitute labels, encode with shifts and masks, and write four bytes.

\begin{lstlisting}[style=cpp]
void outWord(uint32_t w) {
    putchar((w >> 24) & 0xff);
    putchar((w >> 16) & 0xff);
    putchar((w >> 8) & 0xff);
    putchar(w & 0xff);
}
// branch: label resolved relative to PC + 4
int32_t off = ((int32_t)symtab.at(label) - (int32_t)(pc + 4)) / 4;
if (off < -32768 || off > 32767) error("branch out of range");
outWord((op << 26) | (s << 21) | (t << 16) | (off & 0xffff));
\end{lstlisting}

\begin{example}
Addresses are shown in the comments; the arrows give the encoded values.
\begin{lstlisting}[style=mips]
main:  lis $2          ; 0x00
       .word end       ; 0x04  -> 0x00000014
loop:  beq $1, $0, end ; 0x08  -> offset (0x14 - 0x0c)/4 = 2
       sub $1, $1, $2  ; 0x0c
       beq $0, $0, loop; 0x10  -> offset (0x08 - 0x14)/4 = -3
end:   jr $31          ; 0x14
\end{lstlisting}
Symbol table: \code{main}$\mapsto$\code{0x00}, \code{loop}$\mapsto$\code{0x08},
\code{end}$\mapsto$\code{0x14}.
\end{example}

\section{Error checking}

A good assembler rejects every malformed program. Typical checks:

\begin{center}
\begin{tabular}{lp{9cm}}
\toprule
category & examples \\
\midrule
syntax & wrong number of operands; missing commas/parentheses; unknown mnemonic; a
         label after an instruction on the same line \\
registers & register number outside $0$--$31$ \\
immediates & \code{lw}/\code{sw}/\code{beq}/\code{bne} decimal immediate outside
             $[-32768, 32767]$; hex immediate above \code{0xffff} \\
\code{.word} & decimal outside $[-2^{31}, 2^{32}-1]$; hex above \code{0xffffffff} \\
labels & duplicate definition; use of an undefined label; computed branch offset out of
         16-bit range \\
\bottomrule
\end{tabular}
\end{center}

\begin{important}[Ranges]
A \code{.word} accepts any value that fits in 32 bits under \emph{either} interpretation,
so $[-2^{31}, 2^{32}-1]$. A 16-bit immediate in decimal must be in $[-2^{15}, 2^{15}-1]$, but a
hex immediate may be up to \code{0xffff} (it is just a bit pattern). Be careful to parse
large integers into a 64-bit type before range checking, or overflow will hide errors.
\end{important}

Error output convention in CS 241: print a message containing \code{ERROR} to standard
error and stop. A useful structure is to throw an exception from deep inside parsing and
catch it in \code{main}.

\section{Assembler design notes}

\begin{itemize}
  \item Represent each instruction format generically: a table maps mnemonic to
        (format, opcode/funct, operand pattern), e.g.\ \code{add}$\mapsto$(R, 0x20, \code{reg,reg,reg}).
        This avoids eighteen near-identical cases.
  \item Pass 1 can store a parsed intermediate representation, so pass 2 need not re-check syntax.
  \item Running time is $\Theta(n \log n)$ with a balanced-tree map, or expected $\Theta(n)$
        with a hash table, for $n$ lines.
  \item The symbol table is often printed to standard error for debugging (the assignment
        asked for this).
\end{itemize}

\section{Beyond the CS 241 subset}

Real assemblers also support \emph{pseudo-instructions} that expand into several real
instructions (e.g.\ \code{li}, \code{move}, \code{blt}), directives for data and sections
(\code{.data}, \code{.text}, \code{.ascii}), and expressions such as \code{label+8}. Pseudo
instructions complicate pass 1 because one source line may produce a variable number of
words.

\section{Summary}
\begin{itemize}
  \item Forward references force either two passes or backpatching.
  \item Pass 1: walk lines, record label $\mapsto$ address in the symbol table, check syntax,
        advance PC by 4 per instruction or \code{.word}.
  \item Pass 2: encode each instruction; branch offset $= (\mathit{label} - (\mathit{PC}+4))/4$;
        \code{.word label} stores the absolute address.
  \item Check everything: syntax, register numbers, immediate ranges (decimal vs.\ hex),
        duplicate and undefined labels.
\end{itemize}
